% GENERATED FILE. Edit canonical lessons, then run npm run generate:books.
% canonical-source-hash: fnv1a64:fc183045f5ca76a5
% canonical-lessons: TA-C309-vacakar, TA-C309-neyar, TA-C309-talaiyankam, TA-C309-kuru, TA-C309-kelvippatu, TA-R309-review-302-309

\chapter{Readers and Listeners}
\label{ch:ta-readers-and-listeners}
\hlchaptermodality{309}

\begin{chapteropening}
\textbf{By the end of this chapter:} \emph{I can talk about the people who follow the news.}

Complete the last lesson of chapter 309: I can talk about the people who follow the news.
\end{chapteropening}

\section[\texorpdfstring{\ta{வாசகர்} (vācakar)}{வாசகர் (vācakar)}]{\ta{வாசகர்} (vācakar) — a reader}

\label{lesson:TA-C309-vacakar}



\begin{quote}
Before the new one: say the Tamil for to expose, then the Tamil for openly.
\end{quote}

\subsection*{You'll want to know: \ta{வாசகர்}}
\textbf{\ta{வாசகர்}} — \emph{vācakar} — \textquotedblleft{}a reader\textquotedblright{}.

\begin{quote}
\ta{பல} \ta{வாசகர்கள்} \ta{ஆசிரியருக்குக்} \ta{கடிதம்} \ta{எழுதினார்கள்}.
\end{quote}

\begin{quote}
\emph{pala vācakarkaḷ āciriyarukkuk kaṭitam eḻutiṉārkaḷ.}
\end{quote}

\begin{quote}
Many readers wrote letters to the editor.
\end{quote}

Plural \ta{வாசகர்கள்}.

\begin{cousinweb}
\textbf{\ta{வாசகர்}} (\emph{vācakar}) is built on Sanskrit \emph{vācaka}, \textquotedblleft{}speaking, telling\textquotedblright{} (Monier-Williams), with the Tamil -\ta{அர்} of persons.
\end{cousinweb}

\begin{grammarlens}[title={What you've built}]
One more word for telling what happened.
\end{grammarlens}

\subsection*{Your turn}
\begin{itemize}
\raggedright
  \item \emph{Say it:} \emph{vācakar}
  \item \emph{Say it:} \emph{vācakar}, once more
  \item \emph{Say it:} \emph{pala vācakarkaḷ āciriyarukkuk kaṭitam eḻutiṉārkaḷ.} — Many readers wrote letters to the editor.
  \item \emph{Recall:} say the Tamil for to expose, then the Tamil for openly, then say \emph{vācakar} again
\end{itemize}

\begin{tcolorbox}[breakable,colback=teal!4,colframe=teal!35!black,title={Before you move on}]
What does \ta{வாசகர்} mean? (\textquotedblleft{}A reader\textquotedblright{}.) Say it once more.
\end{tcolorbox}
\section[\texorpdfstring{\ta{நேயர்} (nēyar)}{நேயர் (nēyar)}]{\ta{நேயர்} (nēyar) — a listener; a viewer}

\label{lesson:TA-C309-neyar}



\begin{quote}
Before the new one: say the Tamil for openly, then the Tamil for a reader.
\end{quote}

\subsection*{You'll want to know: \ta{நேயர்}}
\textbf{\ta{நேயர்}} — \emph{nēyar} — \textquotedblleft{}a listener; a viewer\textquotedblright{}.

\begin{quote}
\ta{வானொலி} \ta{நேயர்கள்} \ta{பாடல்களை} \ta{விரும்பிக்} \ta{கேட்டார்கள்}.
\end{quote}

\begin{quote}
\emph{vāṉoli nēyarkaḷ pāṭalkaḷai virumpik kēṭṭārkaḷ.}
\end{quote}

\begin{quote}
The radio listeners enjoyed the songs.
\end{quote}

For radio listeners and television viewers: \ta{அன்பான} \ta{நேயர்களே}, dear listeners.

\begin{cousinweb}
\textbf{\ta{நேயர்}} (\emph{nēyar}) is built on \ta{நேயம்}, \textquotedblleft{}affection\textquotedblright{}, usually traced to Sanskrit \emph{snēha}, \textquotedblleft{}love, fondness\textquotedblright{} (CDIAL 13802): a listener is a well-wisher.
\end{cousinweb}

\begin{grammarlens}[title={What you've built}]
One more word for telling what happened.
\end{grammarlens}

\subsection*{Your turn}
\begin{itemize}
\raggedright
  \item \emph{Say it:} \emph{nēyar}
  \item \emph{Say it:} \emph{nēyar}, once more
  \item \emph{Say it:} \emph{vāṉoli nēyarkaḷ pāṭalkaḷai virumpik kēṭṭārkaḷ.} — The radio listeners enjoyed the songs.
  \item \emph{Recall:} say the Tamil for openly, then the Tamil for a reader, then say \emph{nēyar} again
\end{itemize}

\begin{tcolorbox}[breakable,colback=teal!4,colframe=teal!35!black,title={Before you move on}]
What does \ta{நேயர்} mean? (\textquotedblleft{}A listener; a viewer\textquotedblright{}.) Say it once more.
\end{tcolorbox}
\section[\texorpdfstring{\ta{தலையங்கம்} (talaiyaṅkam)}{தலையங்கம் (talaiyaṅkam)}]{\ta{தலையங்கம்} (talaiyaṅkam) — an editorial}

\label{lesson:TA-C309-talaiyankam}



\begin{quote}
Before the new one: say the Tamil for a reader, then the Tamil for a listener.
\end{quote}

\subsection*{You'll want to know: \ta{தலையங்கம்}}
\textbf{\ta{தலையங்கம்}} — \emph{talaiyaṅkam} — \textquotedblleft{}an editorial\textquotedblright{}.

\begin{quote}
\ta{இன்றைய} \ta{தலையங்கம்} \ta{விபத்துகள்} \ta{பற்றியது}.
\end{quote}

\begin{quote}
\emph{iṉṟaiya talaiyaṅkam vipattukaḷ paṟṟiyatu.}
\end{quote}

\begin{quote}
Today's editorial is about accidents.
\end{quote}

\ta{தலை}, head + \ta{அங்கம்}, part: the leading piece.

\begin{cousinweb}
\textbf{\ta{தலையங்கம்}} (\emph{talaiyaṅkam}) joins native \ta{தலை}, \textquotedblleft{}head\textquotedblright{} (DEDR 3103), and \ta{அங்கம்}, \textquotedblleft{}a limb, a part\textquotedblright{}, borrowed from Sanskrit \emph{aṅga} (CDIAL 114).
\end{cousinweb}

\begin{grammarlens}[title={What you've built}]
One more word for telling what happened.
\end{grammarlens}

\subsection*{Your turn}
\begin{itemize}
\raggedright
  \item \emph{Say it:} \emph{talaiyaṅkam}
  \item \emph{Say it:} \emph{talaiyaṅkam}, once more
  \item \emph{Say it:} \emph{iṉṟaiya talaiyaṅkam vipattukaḷ paṟṟiyatu.} — Today's editorial is about accidents.
  \item \emph{Recall:} say the Tamil for a reader, then the Tamil for a listener, then say \emph{talaiyaṅkam} again
\end{itemize}

\begin{tcolorbox}[breakable,colback=teal!4,colframe=teal!35!black,title={Before you move on}]
What does \ta{தலையங்கம்} mean? (\textquotedblleft{}An editorial\textquotedblright{}.) Say it once more.
\end{tcolorbox}
\section[\texorpdfstring{\ta{கூறு} (kūṟu)}{கூறு (kūṟu)}]{\ta{கூறு} (kūṟu) — to state; to tell}

\label{lesson:TA-C309-kuru}



\begin{quote}
Before the new one: say the Tamil for a listener, then the Tamil for an editorial.
\end{quote}

\subsection*{You'll want to know: \ta{கூறு}}
\textbf{\ta{கூறு}} — \emph{kūṟu} — \textquotedblleft{}to state; to tell\textquotedblright{}.

\begin{quote}
\ta{சாட்சி} \ta{நடந்ததை} \ta{எல்லாம்} \ta{கூறினார்}.
\end{quote}

\begin{quote}
\emph{cāṭci naṭantatai ellām kūṟiṉār.}
\end{quote}

\begin{quote}
The witness stated everything that had happened.
\end{quote}

Past (I …): \ta{கூறினேன்}, \emph{kūṟiṉēṉ}.

A written word for \ta{சொல்}; news reports say \ta{கூறினார்}, he stated.

\begin{cousinweb}
\textbf{\ta{கூறு}} (\emph{kūṟu}) is native Dravidian (DEDR 1921): to speak, assert, proclaim.
\end{cousinweb}

\begin{grammarlens}[title={What you've built}]
One more word for telling what happened.
\end{grammarlens}

\subsection*{Your turn}
\begin{itemize}
\raggedright
  \item \emph{Say it:} \emph{kūṟu}
  \item \emph{Say it:} \emph{kūṟu}, once more
  \item \emph{Say it:} \emph{cāṭci naṭantatai ellām kūṟiṉār.} — The witness stated everything that had happened.
  \item \emph{Recall:} say the Tamil for a listener, then the Tamil for an editorial, then say \emph{kūṟu} again
\end{itemize}

\begin{tcolorbox}[breakable,colback=teal!4,colframe=teal!35!black,title={Before you move on}]
What does \ta{கூறு} mean? (\textquotedblleft{}To state; to tell\textquotedblright{}.) Say it once more.
\end{tcolorbox}
\section[\texorpdfstring{\ta{கேள்விப்படு} (kēḷvippaṭu)}{கேள்விப்படு (kēḷvippaṭu)}]{\ta{கேள்விப்படு} (kēḷvippaṭu) — to hear of}

\label{lesson:TA-C309-kelvippatu}



\begin{quote}
Before the new one: say the Tamil for an editorial, then the Tamil for to state.
\end{quote}

\subsection*{You'll want to know: \ta{கேள்விப்படு}}
\textbf{\ta{கேள்விப்படு}} — \emph{kēḷvippaṭu} — \textquotedblleft{}to hear of\textquotedblright{}.

\begin{quote}
\ta{நான்} \ta{அந்தச்} \ta{செய்தியை} \ta{நேற்றுதான்} \ta{கேள்விப்பட்டேன்}.
\end{quote}

\begin{quote}
\emph{nāṉ antac ceytiyai nēṟṟutāṉ kēḷvippaṭṭēṉ.}
\end{quote}

\begin{quote}
I heard of that news only yesterday.
\end{quote}

Past (I …): \ta{கேள்விப்பட்டேன்}, \emph{kēḷvippaṭṭēṉ}.

What you hear about takes -\ta{ஐ}; \ta{கேள்வி} here is hearing, not a question.

\begin{cousinweb}
\textbf{\ta{கேள்விப்படு}} (\emph{kēḷvippaṭu}) is built from \ta{கேள்வி}, \textquotedblleft{}hearing, report\textquotedblright{}, from \ta{கேள்}, \textquotedblleft{}to hear\textquotedblright{} (DEDR 2017), with \ta{படு}, \textquotedblleft{}to undergo\textquotedblright{}.
\end{cousinweb}

\begin{grammarlens}[title={What you've built}]
One more word for telling what happened.
\end{grammarlens}

\subsection*{Your turn}
\begin{itemize}
\raggedright
  \item \emph{Say it:} \emph{kēḷvippaṭu}
  \item \emph{Say it:} \emph{kēḷvippaṭu}, once more
  \item \emph{Say it:} \emph{nāṉ antac ceytiyai nēṟṟutāṉ kēḷvippaṭṭēṉ.} — I heard of that news only yesterday.
  \item \emph{Recall:} say the Tamil for an editorial, then the Tamil for to state, then say \emph{kēḷvippaṭu} again
\end{itemize}

\begin{tcolorbox}[breakable,colback=teal!4,colframe=teal!35!black,title={Before you move on}]
What does \ta{கேள்விப்படு} mean? (\textquotedblleft{}To hear of\textquotedblright{}.) Say it once more.
\end{tcolorbox}
\section[(dialogue)]{Looking back: from After the Disaster to Readers and Listeners}

\label{lesson:TA-R309-review-302-309}



\begin{quote}
Say the Tamil for a disaster and a loss of life.
\end{quote}

\subsection*{Your turn}
Say these aloud:

\begin{itemize}
\raggedright
  \item a disaster, a loss of life, relief aid, to rise again, to ravage
  \item an uproar, a commotion, a stampede, refuge, to drive away
  \item to get flustered, to snatch, to smash, to shake up, to ruin
  \item a reporter, a live broadcast, to broadcast, a presenter, a cameraman
  \item a headline, the front page, an issue, an edition, to bring out
  \item a press release, a spokesperson, a magazine, a titbit, to announce
  \item a sensation, a rumour, to gossip, to expose, openly
  \item a reader, a listener, an editorial, to state, to hear of
\end{itemize}

\begin{tcolorbox}[breakable,colback=teal!4,colframe=teal!35!black,title={Before you move on}]
A disaster? (\textbf{\ta{பேரழிவு}}, \emph{pēraḻivu}.) A loss of life? (\textbf{\ta{உயிரிழப்பு}}, \emph{uyiriḻappu}.) Relief aid? (\textbf{\ta{நிவாரணம்}}, \emph{nivāraṇam}.) To rise again? (\textbf{\ta{மீண்டெழு}}, \emph{mīṇṭeḻu}.) To ravage? (\textbf{\ta{சூறையாடு}}, \emph{cūṟaiyāṭu}.) An uproar? (\textbf{\ta{கூச்சல்}}, \emph{kūccal}.) A commotion? (\textbf{\ta{அமளி}}, \emph{amaḷi}.) A stampede? (\textbf{\ta{கூட்டநெரிசல்}}, \emph{kūṭṭanerical}.) Refuge? (\textbf{\ta{தஞ்சம்}}, \emph{tañcam}.) To drive away? (\textbf{\ta{விரட்டு}}, \emph{viraṭṭu}.) To get flustered? (\textbf{\ta{பதறு}}, \emph{pataṟu}.) To snatch? (\textbf{\ta{பிடுங்கு}}, \emph{piṭuṅku}.) To smash? (\textbf{\ta{தகர்}}, \emph{takar}.) To shake up? (\textbf{\ta{உலுக்கு}}, \emph{ulukku}.) To ruin? (\textbf{\ta{சிதை}}, \emph{citai}.) A reporter? (\textbf{\ta{நிருபர்}}, \emph{nirupar}.) A live broadcast? (\textbf{\ta{நேரலை}}, \emph{nēralai}.) To broadcast? (\textbf{\ta{ஒளிபரப்பு}}, \emph{oḷiparappu}.) A presenter? (\textbf{\ta{தொகுப்பாளர்}}, \emph{tokuppāḷar}.) A cameraman? (\textbf{\ta{ஒளிப்பதிவாளர்}}, \emph{oḷippativāḷar}.) A headline? (\textbf{\ta{தலைப்புச்} \ta{செய்தி}}, \emph{talaippuc ceyti}.) The front page? (\textbf{\ta{முதல்} \ta{பக்கம்}}, \emph{mutal pakkam}.) An issue? (\textbf{\ta{இதழ்}}, \emph{itaḻ}.) An edition? (\textbf{\ta{பதிப்பு}}, \emph{patippu}.) To bring out? (\textbf{\ta{வெளியிடு}}, \emph{veḷiyiṭu}.) A press release? (\textbf{\ta{செய்திக்குறிப்பு}}, \emph{ceytikkuṟippu}.) A spokesperson? (\textbf{\ta{செய்தித்} \ta{தொடர்பாளர்}}, \emph{ceytit toṭarpāḷar}.) A magazine? (\textbf{\ta{பத்திரிகை}}, \emph{pattirikai}.) A titbit? (\textbf{\ta{துணுக்கு}}, \emph{tuṇukku}.) To announce? (\textbf{\ta{அறிவி}}, \emph{aṟivi}.) A sensation? (\textbf{\ta{பரபரப்பு}}, \emph{paraparappu}.) A rumour? (\textbf{\ta{வதந்தி}}, \emph{vatanti}.) To gossip? (\textbf{\ta{கிசுகிசு}}, \emph{kicukicu}.) To expose? (\textbf{\ta{அம்பலப்படுத்து}}, \emph{ampalappaṭuttu}.) Openly? (\textbf{\ta{பகிரங்கமாக}}, \emph{pakiraṅkamāka}.) A reader? (\textbf{\ta{வாசகர்}}, \emph{vācakar}.) A listener? (\textbf{\ta{நேயர்}}, \emph{nēyar}.) An editorial? (\textbf{\ta{தலையங்கம்}}, \emph{talaiyaṅkam}.) To state? (\textbf{\ta{கூறு}}, \emph{kūṟu}.) To hear of? (\textbf{\ta{கேள்விப்படு}}, \emph{kēḷvippaṭu}.)
\end{tcolorbox}
